\section{The decision problem and the scope of completion}\label{sec:introduction}

Optimization-based bound tightening (OBBT) solves auxiliary optimization
problems to improve variable bounds. The improvements can strengthen a
nonlinear relaxation, simplify propagation, and reduce the search tree.
They also cost time. A smaller tree is useful only if the saved work exceeds
the cost of obtaining it. The question studied here is therefore when to
start tightening, when to stop, and when a change in the search state makes
it worth considering again.

The repository's earlier study developed local contraction and stalling
theory and tested iterated OBBT as external preprocessing. Its adaptive
rules already stopped when observed contraction became weak. They did not
produce a net general solver benefit: paying for the first round and
restarting the solver could outweigh later reductions. Those outcomes are
retained in
\href{run:../../research-20260922/iterated-obbt/experiment-report.md}{the September experiment report}.
This continuation does not relabel an observed contraction ratio as a new
stopping theorem. It develops finite certificates for a specified
tightening operator and evaluates a limited policy inside the search.

Three questions must remain separate. First, can a mathematical certificate
bound the possible reduction of a coordinate? Second, can a particular
implementation cheaply obtain such a certificate at a useful search
state? Third, does acting on this information reduce total solving time?
A proof answering the first question does not answer the other two.
The same distinction applies to the word ``complete.'' This report
completes a specified research program by giving explicit mathematical
contracts, proofs, implementations, comparative evidence, and limitations.
It does not resolve every open question about allocation of solver effort.

The central distinction is between a witness feasible in the current
relaxation and a witness that remains available after relaxation rebuilding.
The former bounds the benefit of a frozen round. The latter can protect an
inner box throughout all subsequent rounds under precise monotonicity and
cutoff assumptions. This difference determines what a solver may safely
infer from a failed attempt to tighten a bound.

The report is self-contained at the level of the proved statements.
Section~\ref{sec:foundations} specifies the operator and its scope.
Section~\ref{sec:certificates} proves the finite certificates and their
reuse conditions. Section~\ref{sec:constraints} treats constrained
contraction and its assumptions. Section~\ref{sec:allocation} separates
enhancement cost guarantees from guarantees relative to an independent
baseline. Sections~\ref{sec:implementation} and~\ref{sec:evidence}
describe the actual software and outcomes. The related-work account
in Section~\ref{sec:literature} distinguishes classical ingredients from
the present synthesis. Section~\ref{sec:frontier} records the remaining
frontier without treating an unsuccessful performance hypothesis as an
implementation defect that must be hidden.
